
\documentclass{report}

\input{preamble}
\input{macros}
\input{letterfonts}

\title{\Huge{Storia Moderna e Attualità}\\Sistemi Informativi Aziendali}
\author{\huge{Caberlotto Giovanni}}
\date{}

\begin{document}

\maketitle
\newpage% or \cleardoublepage
% \pdfbookmark[<level>]{<title>}{<dest>}
\pdfbookmark[section]{\contentsname}{toc}
\tableofcontents
\pagebreak

\chapter{La seconda rivoluzione industriale} % (fold)
\label{chap:La seconda rivoluzione industriale}

\section{Dalla prima alla seconda rivoluzione industriale}
\subsection{Le principali caratteristiche}
Nel corsodell'Ottocento la rivoluzione industriale
si diffuse in Europa e negli Stati Uniti, Dal 1870
iniziò la seconda rivoluzione industriale, caratterizzata da:
\begin{itemize}
  \item l'utilizzo della ricerca scientifico nell'innovazione tecnologica;
  \item nuove fonti di energia (elettricità e combustione a petrolio);
  \item lo sviluppo di nuovi settori (chimica, siderurgia, industria automobilistica, ecc.);
  \item monopoli e oligopoli;
  \item la catena di montaggio e la produzione in serie;
  \item la società di massa: tutti iniziarono a fruire degli stessi prodotti (creati in serie) e servizi;
  \item uno Stato sempre più presente nel sistema economico ed evoluto in senso democratico.
\end{itemize}

Questo processo giungerà a maturità, nei paesi più progrediti (Stati Uniti, Gran Bretagna), negli anni Venti e Trenta del Novecento.

\subsection{La funzione della scienza}
A partire dalla seconda metà dell'Ottocento la scienza si legò definitivamente alla tecnica. Molti laboratori universitari furono aperti alle industrie e si affermò la figura professionale dell'ingegniere

\section{La rivoluzione della luce e dei mezzi di comunicazione}
\subsection{L'energia elettrica}
Grazie all'utilizzo dell'energia elettrica e all'invenzione della lampadina a opera di Edison fu possibile diffondere l'illuminazione nelle case, nelle fabbriche e nelle strade: il modo di vivere e di produrre cambio radicalmente

\subsection{Il telegrafo}
Il telegrafo prima e il telefono poi rivoluzionarono la trasmissione dei messaggi, rendendo possibile comunicare anche a grande distanza e in tempi brevissimi. Grazie alle reti di cavi telegrafici, alla fine dell'Ottocento l'Europa era collegata con tutti i continenti.

\subsection{La fotografia e il cinematografo}
All'invenzione del cinematografo segui quella della fotografia. resa possibile dalle nuove conoscenze scientifiche. Nel 1861 comparve la prima macchina fotografica a colori grazie agli esperimenti di J.C. Maxwell. Messo a punto dai fratelli Lumière nel 1895, il cinematografo permise la visione di immagini in movimento. Inizialmente limitato alla riproduzione di scene quotidiane, divenne la fama d'arte tipica del Novecento, conoscendo un enorme successo.

\section{La catena di montaggio e la rivoluzione dei trasporti}
\subsection{Il taylorismo}
Le fabbriche si rinnovarono non solo per l'utilizzo delle grandi innovazioni tecnologiche, ma anche perché il sistema produttivo fu riorganizzato in modo da massimizzare la produzione. Il primo a occuparsi di questo tema fu l'ingegnere americano Taylor. Da lui prende il nome il taylorismo, o organizzazione scientifica del lavoro.

\subsection{Scomporre il processo di produzione}
Secondo Taylor era necessario scomporre il più possibile il processo di produzione di un determinato oggetto. Ciò permetteva di:
\begin{itemize}
  \item affidare a ogni operaio una mansione da ripetere in tempi sempre uguali;
  \item organizzare la fabbrica secondo criteri di efficienza produttiva;
  \item legare i salari degli operai agli effettivi risultati ottenuti (lavoroa cottimo).
\end{itemize}

Il costo della manodopera sarebbe diminuito, i salari sarebbero aumentati e la produzione sarebbe cresciuta. La teoria di Taylor venne applicata per la prima volta su vasta scala da Ford nella sua nota fabbrica di automobili. La catena di montaggio ridusse enormemente i tempi di lavoro, ma lo rese spersonalizzato e ripetitivo.

\subsection{Il successo delle officine ford}
L'intuizione di "portare il lavoro agli uomini anziché gli uomini al lavoro" rivoluzionò il modo di produrre le auto. La Ford T, prodotta alla catena di montaggio e quindi a basso costo divenne l'auto più diffusa e venne prodotta per circa vent'anni in 15 milioni di esemplari.

\subsection{La conquista del cielo: l'aeroplano}
Nel 1903 i fratelk Wright riuscirono a far volare una macchina più pesante dell'aria. Nel 1909 il francese Louis Blériot attraversò la Manica a bordo di un aereo di sua costruzione. nel contempo si sviluppò l'aviazione militare. Da allora la superiorità militare non si sarebbe più conquistata con combattimenti terrestri: l'avrebbe avuta chi avesse dominato i cieli.

\subsection{Il piacere della velocità}
L'automobile e l'aeroplano segnarono l'avvento dell'epoca della velocità nei trasporti. Si diffuse cosi un nuovo modo di intendere lo spazio e il tempo. improntato alla libertà e all'individualismo.

\section{Il capitalismo monopolistico e finanziario}
\subsection{La grande depressione}
Tra il 1870 e il l914 la produzione industriale mondiale quadruplicò, ma negli anni 1873-94 si verificò quella che viene definita la grande depressione. Fu una crisi dovuta alla sovrapproduzione (si produsse in eccesso rispetto alla domanda) a causa:

\begin{itemize}
  \item dell'accresciuta concorrenza internazionale favorita dallosviluppodelle reti di trasporto ferroviario e navale;
  \item dell'incremento produttivo dovuto ai nuovi criteri di produzione, a cui non corrispose una crescita dei salari e della domanda.
\end{itemize}

Ovunque (salvo in Inghilterra) il libero scambio fu sostituito con il proteziordsmo. Le industrie meno competitive fallirono. Quelle che si riorganizzarono uscirono dalla crisi rafforzate e si ingrandirono. Le banche concessero prestiti alle aziende importanti piuttosto che a quelle piccole, ma per stare sul mercato occorrevano investimenti massicci.Così poche imprese assunsero il controllo del mercato: nacque il fenomeno della concentrazione industriale.

\subsection{Un nuovo tipo di capitalismo}
La concentrazione industriale assunse varie forme, tra cui il monopolio. che si verifica quando un'unica impresa controlla un settore produttivo. Se le imprese sono poche, si ha invece un oligopolio, Il capitalismo oltre che monopolistico, divenne finanziario in quanto l'interesse finanziario prese a dominare su quello industriale. Di conseguenza crebbe l'importanza delle banche che concedevano prestiti alle industrie.

\section{La crescità demografica e la nascità della medicina moderna}

\subsection{Il boom demografico}
Tra il 1850 e il 1914 vi fu un boom demografico nei paesi più arretrati. Nei paesi industrializzati, nonostante il miglioramento delle condizioni di vita, la natalità prese a decrescere per:
\begin{itemize}
  \item l'innalzamento della scolarità;
  \item l'inserimento delle donne nel sistema produttivo;
  \item la diffusione dei metodi di controllo delle nascite.
\end{itemize}

\subsection{La scoperta dell'igene}
Lo studio e la cura delle malattie infettive ebbero grande impulso dall'osservazione degli ambienti dove esse si sviluppavano. La carenza di igiene di abitazioni e interi quartieri favoriva la comparsa e la diffusione di terribili malattie: diversi medici e scienziati imposero alle autorità l'adozione di misure igieniche anche nella costruzione delle case e il risanamento dei quartieri più poveri. L'imposizione di pratiche igieniche anche negli ospedali fece crollare il numero delle infezioni che anche in quei luoghi si sviluppavano.

\subsection{II progresso della medicina}
I grandi progressi nell'osservazione al microscopio portarono medici come Louis Pasteure Robert Koch a isolare i microrganismi responsabili di malattie come la rabbia, la peste e la tubercolosi, e a preparare vaceini per la loro prevenzione o farmaci per la loro cura. Negli ospedali si diffuse la pratica dell'anestesia che rese più semplici le operazioni.

\subsection{La diffusione delle droghe}
Oltre alle droghe già conosciute come l'oppio e l'hashish le più approfondite conoscenze chimiche portarono alla nascita di nuove sostanze come la morfina,un antidolorifico che però causava dipendenza psichica e fisica. Anche l'eroina nacque come farmaco antidolorifico. Di fronte alle conseguenze negative determinate dall'uso dell'eroina, l'industria Bayer che
l'aveva prodotta la ritirò dal commercio. Al suo posto venne lanciato un farmaco destinato ad avere un immenso e positivo successo: l'aspirina. Nel contempo però si diffondeva anche un'altra nuova droga, la cocaina,un alcaloide estratto per la prima volta nel 1860 dalle foglie della pianta della coca.

\subsection{L'emigrazione}
I paesi industrializzati reagirono alla crisi agricola ammodernando il settore, mentre in quelli arretrati soprattutto dell'Europa centroorientale si creò un esubero di popolazione nelle campagne.Iniziarono così i flussi migratori verso le città, verso altri paesi d'Europa e oltreoceano. Il movimento migratorio dall'Europa verso gli Stati Uniti che si manifestò alla fine del XIX secolo rappresenta un fatto senza precedenti, e mai più ripetuto con le medesime proporzioni. Per comprenderne l'entità, basti pensare che dei 55 milioni di Europei emigrati dal 1821 al 1924, la maggior parte, 21 milioni, partirono tra il 1870 e il 1900.

% chapter La seconda rivoluzione industriale (end)

\chapter{La spartizione imperialistica del mondo} % (fold)
\label{chap:La spartizione imperialistica del mondo}

\section{L'imperialismo: la competizione globale}
\subsection{Definizione}
L'imperialismo fu una corsa alla colonializzazione guidata da governi in competizioni tra loro, che ebbe come obbiettivo l'estensione dei confini nazionali. Tra il 1870 e il 1914 un quarto del mondo venne spartito fra pochi stati.

\subsection{Il contesto politico}
La Germania era diventata il punto di equilibrio dei rapporti di forza in Europa. Bismarck aveva garantito la pace con una politica di equilibrio, ma la tensione salì a causa:
\begin{itemize}
  \item del revanscismo francese;
  \item delle tensioni nei Balcani;
  \item della competizione coloniale.
\end{itemize}

\subsection{Il contesto economico e culturale}
Con la «grande depressione» (1873-96) gli Stati presero a sostenere l'economia nazionale con il protezionismo, le commesse stata e la politica imperialista. Quest'ultima:
\begin{itemize}
  \item garanti nuovi sbocchi commerciali e materie
prime a basso costo;
  \item fu sorretta da motivazioni ideologiche, fondate sul nazionalismo, sul razzismo e sul mito della missione civilizzatrice degli Europei: il «fardello desinane bianco» (Kipling).
\end{itemize}

\section{La spartizione dell'africa e la conferenza di berlino}
\subsection{L'apertura del canale di Suez}
L'apertura di un canale fra il Mediterraneo e il Mar Rosso, tagliando l'istmo diSuez, fu fondamentale per la spinta alla colonizzazione dell'Africa. Organizzata dai Francesi, l'opera richiese 10 anni di lavoro e circa 400 miliardi di franchi. Il commercio internazionale conobbe una svolta: il tempo impiegato peri viaggi in Oriente veniva in pratica dimezzato. Nel 1882 l'Inghilterra ottenne il controllo del Canale.

\subsection{L'espansione in Africa}
La Francia, che possedeva già l'Algeria, intendeva estendere il suo impero coloniale lungo l'asse ovest-est dell'Africa centrosettentrionale partendo dalla Tunisia ccupata nel 1881.
L'espansione dell'Inghilterra invece si spingeva lungo l'asse nord—sud partendo dall'Egitto occupato nel 1882.

\subsection{La conferenza di Berlino}
La conferenza di Berlino (1884-85) sancì il principio dell'occupazione di fatto come criterio di possesso dei territori africani: ciò scateno ancora di più la competizione coloniale, con il coinvolgimento della stessa Germania, ultima arrivata nella corsa coloniale.

\subsection{Boeri e Inglesi}
Nella zona Sudafricana la scoperta dei giacimenti d'oro e di diamanti scatenò il conflitto tra i Boeri (vecchi coloni Olandesi) e Inglesi quest'ultimi ebbero la meglio e nel 1910 formarono l'unione Sudafricana.


\section{La spartizione dell'Asia e l'espansionismo americano}
\subsection{La colonizzazione in Asia}
La possibilità di accedere dal Mar Mediterraneo al Mar Rosso tramite il canale di Suez, costruito tra il 1859 e il 1869, diede un nuovo impulso all'espansione europea in Asia, già in atto prima dell'età dell'imperialismo.


\subsection{Il dominio Inglese in India}
Dopo lo scoppio della rivolta dei sepoy (soldati indiani arruolati nell'esercito inglese), nel 1857 il governo britannico assunse il controllo diretto dell India. prima governata per conto della Gran Bretagna dalla Compagnia delle Indie. Nel 1885 nacque il Consesso Nazionale Indiano un'assemblea della classe media indiana che passò da posizioni di collaborazione con gli Inglesi alla richiesta di autogoverno.

\subsection{La Cina}
Le tensioni tra gli Inglesi e il governo imperiale causarono due guerre dell'oppio (tra il 1839 e il 1860 così chiamate perché determinate dal rifiuto cinese di importare l'oppio Le vare potenze approfittarono della crisi dell'Impero cinese per occuparne alcune zone. Una rivolta condotta dalla società segreta xenofoba dei Boxer venne sedata da un contingente internaiionale: ormai la sovranità del governo imperiale era limitata.

\subsection{L'espansionismo americano}
A fine Ottocento gli Stati Uniti abbandonarono l'isolazionismo ed entrarono nella competizione imperialista. puntando non alla conquista territoriale ma all'egemonia economica. Il loro primo obiettivo fu l'America Latina.

% chapter La spartizione imperialistica del mondo (end)


\chapter{Dalla società di massa alla globalizzazione} % (fold)
\label{chap:Dalla società di massa alla globalizzazione}

\section{Che cos'è la società di massa}

\subsection{Definizione}
che cosa intendamo con società di massa? La risposta più immediata può essere: la società attuale.
Il concetto di "società di massa" si riferisce a una società caratterizzata dalla presenza diffusa di determinate caratteristiche sociali, economiche e culturali. Questa idea è stata ampiamente esplorata e discussa dai sociologi e dagli studiosi delle scienze sociali nel XX secolo.
Nelle società di massa, si osservano diversi tratti distintivi:

\begin{enumerate}
  \item Popolazione numerosa: Le società di massa hanno una grande popolazione in cui gli individui spesso non si conoscono personalmente. Ciò può portare a una mancanza di legami sociali intimi e a una maggiore anonimicità tra le persone;
  \item Standardizazione: Le attività quotidiane, come l'istruzione, il lavoro e l'intrattenimento, tendono a essere standardizzate e omogenee per soddisfare le esigenze di un gran numero di persone. Questo può portare alla perdita di specificità culturale e di identità individuale;
  \item Media di massa: Le informazioni e l'intrattenimento sono diffusi attraverso i media di massa, come la televisione, la radio e, più recentemente, internet. Questi mezzi hanno un vasto pubblico e possono influenzare le opinioni e i comportamenti di molte persone;
  \item Consumo di massa: Le società di massa vedono un aumento della produzione e del consumo di beni e servizi su larga scala. Ciò può portare a un'economia orientata al consumo, in cui le persone sono spinte a acquistare prodotti di massa;
  \item Urbanizzazione: Le società di massa spesso presentano alti livelli di urbanizzazione, con un gran numero di persone che vivono in città e aree metropolitane. Questo può portare a sfide come l'anonimato sociale e la perdita di legami comunitari;
  \item Mobilità sociale: Nelle società di massa, l'ascesa sociale è spesso basata sul merito e sull'istruzione, consentendo alle persone di spostarsi socialmente in base alle loro capacità e alle loro realizzazioni personali.
\end{enumerate}

È importante notare che il concetto di società di massa è stato oggetto di critica e dibattito. Alcuni sostengono che questa teoria semplifichi eccessivamente la complessità delle società moderne, mentre altri la vedono come uno strumento utile per comprendere alcune delle tendenze sociali e culturali nelle società industrializzate e globalizzate.

\subsection{L'espansione del terziario}
Nell'economia delle società di massa assunse un ruolo sempre più importante il terziario, ovvero il settore economico che riguarda i servizi: commercio, banche, ospedali, scuole ecc. 
Furono i lavoratori del terziario, sia dipendenti sia autonomi, a ingrossare le file del ceto medio e ad accrescerne il peso sulla società.
Il progressivo allargamento delle competenze dello Stato portò all'aumento dei dipendenti pubblici, impiegati nei trasporti, nella sanità nell'istruzione ecc. Ma crebbe anche il numero dei dipendenti di ditte private (impiegati, commessi, tecnici specialisti) che, a differenza degli operai, svolgevano mansioni non manuali. Il ceto medio faceva propri i valori tradizionali della borghesia: indvidualismo, rispettabilità, proprietà privata, deferenza nei confronti delle gerarchie, patriottismo.

\subsection{Partiti di massa}
Con l'introduzione del suffraggio universale maschile nella maggiorparte dell'territorio europeo si notò un cambiamento a livello politico, per conquistare il consenso di masse di elettori sempre in aumento si affermarono progressivamente i partiti politici di massa. Contemporaneamente, sorsero organizzazioni sindacali non più ristrette e territoriali ma a livello nazionale.

\subsection{La vita quotidiana}
La diffusione su larga scala dei beni di consumo rese più comoda la vita di molte famiglie, ma il cambiamento, rispetto al passato, non riguardò solo l'utilizzo di beni e comodità materiali. Insieme all'illuminazione elettrica, all'acqua potabile e alle automobili si diffusero anche i mass media (media di massa) sopra citati. Fu sopratutto la stampa quotidiana e periodica a incrementare la produzione e la vendita. Anche il mondo della scuola subì un mutamento sostanziale: l'istruzione non venne più considerata un bene elitario, riservato a nobili e abbienti, ma un'opportunità da offrire a tutti i cittadini. La posta in gioco era superare quella che ormai si considerava una piaga sociale: l'analfabetismo

\subsection{La moda: liberarsi dal superfluo}
Agli inizi del Novecento, l'espansione del mercato mondiale permise di utilizzare diversi tipi di tessuti preziosi e colorati, resistenti e sofisticati. Anche i gusti mutarono, rivolgendosi verso motivi esotici. Si affermò difatti la tendenza a privilegiare i colori vivaci, i disegni orientaleggianti, le decorazioni. La nuova tecnologia applicata al tessile portò sul mercato nuove fibre artificiali, come il rayon, derivato dalla cellulosa. La distribuzione di capi di vestiario iniziò a diffondersi anche nei grandi magazzini, dove i prezzi potevano essere più convenienti. La nuova moda non stravolgeva più il corpo della donna, ma lo liberava. L'imperativo sembrava essere quello di liberarsi del superfluo.

\subsection{Giochi di massa: le olimpiadi moderne}
Nell'Ottocento lo sport e la ginnastica divennero un importante momento educativo.
Fu l'inghilterra a inserire per prima la ginnastica tra le materie della scuola dell'obbligo. Ma la vera rinascita dello sport agonistico come momento di confronto fisico e spirituale si deve a Pierre de Coubertin, che propose di svolgere di nuovo i Giochi olimpici, gloria dell'antichità greca. Venne scelta come sede Atene e le prime olimpiadi moderne si svolsero nel 1896. Tuttavia le olimpiadi moderne erano ben diverse da quelle del passato: non più un confronto fisico e spirituale, ma la vetrina sportiva della società di massa

\section{Il dibattito politico e sociale}

\subsection{L'eredità dell'ottocento}
Il dibattito ideologico dell'ottocento si era sviluppato sui grandi problemi causati dalla rivoluzione industriale e sopratutto sulla questione sociale.
I conservatori chiedevano allo Stato di reprimere tutte le agitazioni popolari.
I liberali chiedevano allo Stato di non intervenire nell'economia. 
Le uniche leggi che dovevano regolare il mercato erano quelle della libera concorrenza.
I socialisti sostenevano invece che una società più giusta non potesse nascere se non dalle lotte dei più oppressi: agricoltori e operai.
La chiesa condannava sia il socialismo che il libero mercato, invitando imprenditori e lavoratori ad abbandonare lo scontro e a realizzare una collaborazione pacifica.

\subsection{Il socialismo in europa}
Nel corso dell'ottocento all'interno del movimento socialista si impose la tendenza marxista, che individuava nella rivoluzione lo strumento del riscatto del ploletariato. Sul finire del secolo in tutti i paesi europei sorsero i partiti socialisti. Il primo a formarsi, e il più importante fu l'SPD, il Partito Socialdemocratico Tedesco, nato nel 1875. In Italia il partito socialista venne fondato nel 1892. Solo in Gran Bretagna il marxismo non riuscì a imporsi all'interno dei sindacati, dove per iniziativa degli stessi dirigenti sindacali sorse il partito laburista

\subsection{La seconda internazionale}
Nel 1899, per costituire la Seconda internazionale socialista, tutti i partiti socialisti europei si riunirono a Parigi, dove vennero approvate la limitazione delle giornata lavorativa a otto ore e la proclamazione di una giornata mondiale di lotta per il primo maggio di ogni anno. 
A dominare la seconda internazionale fu il marxismo. tuttavia all'interno dei partiti socialisti si delinearono due tendenze:
\begin{itemize}
  \item Quella revisionista o socialdemocratica, che "rivedeva" i fondamenti stessi dell'analisi marxista, rifiutando la rivoluzione e ponendo in luce la necessità di un'azione democratica e riformista;
  \item quella ortodossa o rivoluzionaria, che non rinunciava a ricorrere alla rivoluzione per raggiungere una società senza classi.
\end{itemize}

\subsection{Sorel e il sindacalismo rivoluzionario}
In Francia ebbe origine un'altra tendenza all'interno del movimento operaio, che prese il nome di sindacalismo rivoluzionario. I sindacalisti francesi insistevano sulla necessità di addestrare le masse operaie alla lotta. Lo sciopero era considerato una "ginnastica rivoluzionaria", in vista del grande sciopero rivoluzionario che avrebbe segnato la fine della società borghese. Georges Sorel fu l'interprete più autorevole di questa tendenza.


\subsection{La dottrina della chiesa cattolica}
Leone XIII, con la pubblicazione della Rerum Novarum, cercò di formulare una proposta sociale coerente con il messaggio evangelico. L'enciclica conteneva le seguenti indicazioni:

\begin{itemize}
  \item Condanna del liberismo perchè privo di preoccupazioni morali in ambito economico;
  \item Condanna delle teorie socialiste e collettivistiche in quanto negavano il "diritto naturale" alla proprietà privata;
  \item Richiesta di un intervento dello Stato per placare il conflitto sociale;
  \item Condanna della lorra di classe ed estrazione a alla collaborazione tra i padroni e gli operai.
\end{itemize}

\subsection{Democrazia cristiana e modernismo}
Negli ultimi anni dell'ottocento emerse nel mondo cattolico una nuova tendenza politica che fu definita "democrazia cristiana".
Secondo i democratici cristiani, bisognava costruire un partito democratico ispirato ai valori del cristianesimo.
Sotto il pontificato di Pio X fu molto attenuato il non expedit, poi abrogato formalmente nel 1919, anno in cui i cattolici poterono dar vita a un loro partito popolare italiano fondato da Luigi Sturzo.
fu sempre Pio X a condannare il modernismo (in quanto incompatibile con la tradizionale dottrina della Chiesa): un movimento, per nulla omogeneo, che si proponeva di reinterpretare la dottrina cattolica in chiave moderna.

\subsection{Suffragette e femministe}
Nella seconda metà dell'ottocento nacquero in Inghilterra e negli Stati Uniti i primi movimenti delle suffragette, che rivendicavano l'estensione del diritto di voto alle donne. Le sufragette erano perlopiù di estrazione sociale borghese; l'obbiettivo del movimento era non solo la possibilità per le donne di votare e di essere elette, ma anche la completa parità tra uomo e donna.

\subsection{Legislazione sociale e sistema fiscale}
Tra la fine dell'Ottocento e gli inizi del Novecento, le pressioni dei partiti di ispirazione socialista e del crescente movimento sindacale spinsero i governi dei maggiori Stati europei a introdurre forme di legislazione sociale sul modello di quelle bismarckiane. Furono istituite assicurazioni contro gli infortuni e di previdenza per la vecchiaia, nonché, in casi limitatim sussidi per i disoccupati. Vennero imposti limiti all'orario di lavoro e al lavoro dei fanciulli.
Queste iniziative fecero aumentare le tasse. La scelta prevalente fu di accresere il peso delle imposte dirette rispetto a quello delle imposte indirette. Venne anche introdotto il principio della tassazzione progressiva: le aliquote di prelievo aumentavano in proporzione al crescere del reddito.
% chapter Dalla società di massa alla globalizzazione (end)

\section{Il nuovo contesto culturale}
\subsection{La crisi del positivismo e l'irrazionalismo}
Tra Ottocento e Novecento, il contesto culturale mutò: il positivismo si rivelò sempre meno adeguato a interpretare i profondi cambiamenti sociali, politici ed economici. Si diffusero così correnti filosofiche irrazionalistiche e vitalistiche, secondo le quali la realtà umana in tutte le sue manifestazioni dipendeva da fattori come l'istinto, l'intuizione e la volontà. Dure critiche al positivismo vennero anche da Friedrich Nietzsche: La visione lineare della storia e del tempo, fondamento dell'idea di progresso, veniva da lui criticata come un'illusoria rilettura della tradizione ebraico-cristiana ormai in decadenza.

\subsection{Idealismo, storicismo e pragmatismo}

In Italia l'opposizione alla mentalità positivista sfociò in una ripresa della filosofia idealistica, i cui principali esponenti furono Benedetto Croce e Giovanni Gentile.
In Germania nacque una corrente chiamata storicismo (il cui principale esponente fu Wilhelm Dilthey) che riconosceva alla Storia un carattere scientifico.
Nei paesi Anglossasoni e in particolare negli Stati Uniti, si sviluppò il pragmatismo: questa corrente filosofica riteneva determinante l'apporto della "pratica"

\subsection{La psicologia e la psicanalisi}
Nacque la moderna psicologia, come studio delle funzioni mentali che stanno alla base del comportamento umano.
Intorno agli anni Settanta del XIX secolo fu il tedesco Wilhelm Wundt ad applicare metodologie sperimentali e quantitative nello studio della psiche.
Grandissima influenza ebbe Sigmund Freud, fondatore della psicanalisi. Secondo Freud, la psiche umana mostra una vita cosiente (l'io), ma anche la presenza di una vita incosciente (l'Es) che appare guidata minacciosamente da leggi proprie. Freud riteneva che i contenuti inconsci causa di nevrosi potessero essere individuati mediante l'interpretazione dei sogni.

\subsection{Il pensiero politico}
Gaetano Mosca formulò la teoria della classe politica, secondo la quale, anche nelle società democratiche a sovranità popolare, il potere effettivo è detenuto da un ristretto numero di individui: i politici di professione, la "classe politica" o "dirigente". per Vilfredo Pareto la politica era dominata da istinti e non era altro che uno scontro di élite (gruppi qualificati, oligarchie). Roberto Michels riconobbe un legame inevitabile fra la nascita dei partiti organizzati di massa e la comparsa di oligarchie burocratiche, vere protagoniste della politica. Secondo Max Weber l'accrescersi degli apparati burocratrici era una tendenza inevitabile in quella fase di sviluppo della società: la sua ramificazione e pervasività mettevano a rischio la libertà dell'individuo.

\subsection{I progressi delle scienze}
Alla fine dell'Ottocento gli Inglesi Joseph Thompson ed Ernest Rutherford posero le basi della disica atomica: il tedesco Max Plank formulò, nel 1900, le basi della fisica quantistica; nel 1905 il tedesco Albert Einstein enunciò la teoria della relatività: dopo queste scoperte il mondo apparve agli scienziati molto diverso. Einstein affermava che la misurazione del tempo poteva variare in funzione di altre variabili, ad esempio la velocità. ciò significava però che, nella grande estensione dell'universo, potevano essere validi modelli diversi. La fiducia ingenua nelle "scienze esatte", punto di forza del positivismo, ne risultava quindi indebolita.

\end{document}
